\scanpage{53}
sur \deleted{\unsure{des}} \inserted{chacun} des \unsure{ens} des quatre sommets
$a_i,\alpha_i,a'_j,a'_k$ \unsure{se} trouvant sur un $D_{a_i}$,
i.e. en \unsure{précisant|disant} lequel, parmi les trois
sommets $\alpha_i,a'_j,a'_h$ distincts de $a_i$, est
\uline{opposé} à $a_i$. On voit que si la
figure \inserted{1-dimensionnelle} provient d'un arrang\textsuperscript{t} 2-dimensionnel,
alors il y a \inserted{(Cas I)} un ordre circulaire
% slanted
\margin{Il y a une autre possibilité, c'est que pour \deleted{\unsure{tout}} $i\in J$, $a_i$ soit opposé à $\alpha_i$ (cas II)}
sur $J$ (unique), tel que $a_i$ soit opposé,
sur $D_{a_i}$, à $a'_{i+1}$. (\uline{NB} Dans le cas
de figure choisi, cet ordre circulaire est
$\unsure{abc^{\circlearrowright}|abc^{\rightarrow}|abc}$).
\deleted{On \unsure{fera|verra|trouve}} \deleted{On voit que les configurations} \deleted{\unsure{Quand}}
\deleted{2-dimensionnelles} La situation $K$, donc
$K_{J'}$, \deleted{\unsure{prov}} et $K_J$, provient d'un arr. 2-dimensionnel,
alors, \deleted{\unsure{les}} désignant par
$\lambda_i$ $(i\in J)$ \deleted{les} \inserted{l'}arête de $K_J$ \deleted{qui con} portée
par $D_i$ qui contient $\alpha_i$, les $\lambda_i$ $(i\in J)$
font un syst. d'arêtes compatible avec
l'orientation dans le cas I, et non compatible
dans le cas II.

(A6) \deleted{L'orien} Dans le cas d'un système
de trois ps droites $J\in\unsure{\mathfrak P_3|\mathcal P_3}(I)$ non concourantes,
tel que \unsure{tt|toute} autre ps. droite passe par un des
trois sommets \inserted{$a_i$} de $K_J$, et lorsqu'on peut trouver
pour chaque $a_i$ une $d_{a_i}$ qui passe par $a_i$ et
% archive sheet number: 52
